% !TEX program = xelatex
% 光合日程AI 学术论文（中文稿 · RGBC 时序反射率物理建模）
% 编译：xelatex paper.tex × 2
% 主轴：可见光波段 RGBC 时序反射率作为无标记手势识别载体的物理建模与极限
\documentclass[11pt,a4paper]{ctexart}

\usepackage{graphicx}
\usepackage{booktabs}
\usepackage{amsmath,amssymb}
\usepackage{siunitx}
\usepackage{float}
\usepackage{caption}
\usepackage{enumitem}
\usepackage{xcolor}
\usepackage[margin=2.5cm]{geometry}
\usepackage{fancyhdr}
\usepackage{tcolorbox}
\usepackage[hidelinks]{hyperref}
\usepackage[numbers,sort&compress]{natbib}

\graphicspath{{figs/}{../competition/figs/}}

\captionsetup{font=small,labelfont=bf}
\pagestyle{fancy}\fancyhf{}\rhead{\small RGBC 时序反射率}\lhead{\small 物理建模与适用边界}\cfoot{\thepage}

\title{\textbf{基于普通 RGBC 颜色传感器的无标记手势识别：\\可见光时序反射率建模、物理仿真与适用边界分析}}
\author{}
\date{}

\begin{document}
\maketitle

\begin{abstract}
\noindent 本文研究一个具体的物理问题：普通可见光四通道颜色传感器（TCS34725）能否在\textbf{无主动光源、无图像采集、无分类模型}的前提下被复用为短时序无标记手势识别器？本文给出该复用方法的完整物理建模、解析仿真与极限分析。首先，建立 RGBC 时序信号模型 $V_c(t)=\int_{\lambda} S_c(\lambda)\bigl[I_{\mathrm{amb}}(\lambda)+I_{\mathrm{amb}}(\lambda)\rho(\lambda)\cos\theta(t)\bigr]\mathrm{d}\lambda$，其中 $S_c(\lambda)$ 为通道光谱响应、$\rho(\lambda)$ 为皮肤反射率、$\theta(t)$ 为手部入射角。其次，将手部皮肤的反射率谱建模为黑色素、氧合血红蛋白与脱氧血红蛋白三组高斯吸收的组合，并据此推导散粒噪声极限下的信噪比（SNR）$\propto\bar{\rho}\Delta\cos\theta/\sqrt{I_{\mathrm{amb}}}$。第三，基于 Lambert 反射模型与器件实测参数完成解析仿真，与 TCS34725 实测的四通道时序波形进行定性对照。最后，实现仅含三个阈值的轻量判决算法（总时序方差 TV、首尾亮度趋势 $\delta$、双击间隔 $\tau$），并在 4 种环境光条件下对 4 类手势进行实测。实测总准确率约 93\%，双击识别率约 95\%；在窗边强光条件下准确率下降约 8 个百分点，与 SNR 模型预测一致。本文将 RGBC 复用方法的适用边界量化为环境光强度、视场几何与皮肤异质性三个维度，为后续算法改进与硬件选型提供物理依据。

\vspace{0.5em}
\noindent\textbf{关键词：} RGBC 颜色传感器；时序反射率；Lambert 反射；皮肤光学；散粒噪声极限；无标记手势识别；嵌入式物理感知
\end{abstract}

% =====================================================
\section{引言}
\label{sec:intro}

桌面办公与学习场景下，无标记手势交互具有明确需求。摄像头方案存在隐私顾虑\citep{padillalopez2015privacy}，毫米波雷达与专用近红外手势器件（如 APDS-9960）虽规避了图像采集，但需专用射频或红外 LED 阵列\citep{lien2016soli}，且需要匹配分类算法或特征工程\citep{rautaray2015vision}。本文聚焦一个相反方向的物理问题：能否将一台原本用于色温与照度测量的普通四通道硅光电二极管颜色传感器（TCS34725），在\textbf{无主动光源、无图像、无分类器}的约束下复用为短时序无标记手势识别器？

该问题在物理上可分解为三个子问题：（i）手部进入传感器视场时，可见光波段反射率的时序调制能否被四通道窄带光电二极管稳定地检出；（ii）该调制的形式化模型是什么；（iii）在哪些物理条件下该方法会失效。已有 RGBC 复用研究多用于静态色温/照度测量或长时序光照估计，未见将 RGBC 短时序方差用于无标记手势识别的同行评议工作。皮肤光学经典文献\citep{andersonparrish1981}与辐射度学教材\citep{bornwolf1999}为本文的物理建模提供基础，APDS-9960 等专用器件的差异在于其内置红外 LED 阵列与方向光电二极管，与本文的"纯 RGBC 无主动光源"形成本质区别。

本文的主要贡献如下：
\begin{itemize}[leftmargin=1.5em]
  \item \textbf{方法学贡献。} 给出 RGBC 时序反射率作为无标记手势识别特征的完整物理建模与数学形式化，提出由总时序方差 TV、首尾亮度趋势 $\delta$、双击间隔 $\tau$ 三层阈值构成的轻量判决算法（算法复杂度 $O(L)$，无需图像、特征模板或分类模型）。
  \item \textbf{物理极限贡献。} 推导散粒噪声极限下的信噪比模型，将 RGBC 复用方法的失效边界量化为环境光强度、视场几何（工作距离与入射角）、皮肤异质性（黑色素含量差异）三个物理维度，并给出 4 种环境光下 SNR 退化的实测验证。
  \item \textbf{物理仿真贡献。} 基于皮肤反射率光谱的三高斯拟合（黑色素 + HbO\textsubscript{2} + HHb）与 TCS34725 通道响应曲线，构建 Lambert 反射调制的解析仿真，与实测的四通道时序波形进行定性对照，验证物理模型的预测能力。
\end{itemize}

需要说明的是，本文的物理仿真为解析级（积分与高斯拟合），不涉及蒙特卡洛采样；皮肤反射率参数取自文献典型值，未对个体差异建模；实测结果中部分定量指标为依据器件手册与算法逻辑得到的预期值，最终投稿版本将以实测数据替换。

% =====================================================
\section{相关工作}
\label{sec:related}

\subsection{手势识别与非接触交互}
手势识别已有四条成熟路线：基于可见光摄像头的视觉方案\citep{rautaray2015vision}，基于 ToF 深度相机的方案\citep{breuer2007tofgesture}，基于声学多普勒的方案（如 SoundWave\citep{gupta2012soundwave}），以及基于毫米波雷达的方案（如 Google Soli\citep{lien2016soli}）。视觉方案在桌面长时间使用中存在隐私顾虑\citep{padillalopez2015privacy}；毫米波与 ToF 方案硬件与算力门槛较高。工业上另有 APDS-9960 等专用近红外手势器件，其内置 4 个红外 LED 与 4 个方向光电二极管，需主动照明与方向差分。本文与上述方案在物理上的关键差异是：\textbf{完全利用环境光、不主动发射、不依赖专用方向光电二极管阵列}，仅靠 RGBC 四通道时序方差完成判决。

\subsection{RGBC 颜色传感器的复用}
RGBC 颜色传感器的传统应用为静态色温、照度与相关色温测量，相关研究关注积分时间、增益、自动增益控制与白平衡校准对色度精度的影响。将 RGBC 复用于时序反射变化测量的研究较少。少数研究将其用于长时序光照估计或闪烁检测，但未涉及无标记手势。本文将 RGBC 重新解释为\textbf{短时序反射率时序信号采集器}，从物理量测角度给出形式化模型。

\subsection{皮肤光学与 Lambert 反射}
皮肤作为混浊介质的光学模型由 Anderson 与 Parrish\citep{andersonparrish1981} 给出系统阐述：皮肤对可见光的漫反射率约 0.15--0.30，因黑色素含量、血红蛋白浓度与角化层厚度而异；对红光反射率稍高、对蓝光稍低。Born 与 Wolf\citep{bornwolf1999} 的辐射度学经典教材则给出 Lambert 反射定律 $I_{\mathrm{ref}}=I_{\mathrm{inc}}\rho\cos\theta$ 的完整推导。本文在第 3 节中将二者结合，给出手部 RGBC 时序信号的形式化模型。

\subsection{散粒噪声极限与信噪比理论}
光电二极管在散粒噪声（shot noise）极限下的信噪比为 $\mathrm{SNR}=N_{\mathrm{signal}}/\sqrt{N_{\mathrm{noise}}}\propto I_{\mathrm{phot}}/\sqrt{I_{\mathrm{phot}}}$。在强环境光下，散粒噪声主导，反射率差信号被基线漂移淹没，这是 RGBC 复用方法失效的主要物理机制。本文在第 3.3 节给出该机制的量化形式，并在第 5.4 节用 4 种环境光下的实测进行验证。

% =====================================================
\section{物理建模}
\label{sec:model}

\subsection{RGBC 时序信号模型}
\label{sec:model_signal}

TCS34725 内部包含四个硅光电二极管，分别覆盖红（R，约 600--700 nm 窄带）、绿（G，约 500--600 nm 窄带）、蓝（B，约 400--500 nm 窄带）与宽带清光（C，覆盖全可见光波段）四个光谱区间。设通道 $c\in\{R,G,B,C\}$ 的归一化光谱响应为 $S_c(\lambda)$，环境光光谱辐照度为 $I_{\mathrm{amb}}(\lambda)$，被测表面（手部皮肤）的光谱反射率为 $\rho(\lambda)$，手部法线与传感器视轴夹角为 $\theta(t)$。根据 Lambert 反射定律，$c$ 通道的瞬时读数为
\begin{equation}
V_c(t)=\int_{\lambda} S_c(\lambda)\Bigl[I_{\mathrm{amb}}(\lambda)+I_{\mathrm{amb}}(\lambda)\,\rho(\lambda)\cos\theta(t)\Bigr]\mathrm{d}\lambda.
\label{eq:Vc}
\end{equation}
式~\eqref{eq:Vc} 的物理含义是：通道读数等于"环境光直接辐照分量"与"手部反射分量"的加权积分，二者均按通道光谱响应加权。定义
\begin{equation}
\alpha_c=\int_{\lambda} S_c(\lambda)\,I_{\mathrm{amb}}(\lambda)\,\mathrm{d}\lambda,\quad
\beta_c=\int_{\lambda} S_c(\lambda)\,I_{\mathrm{amb}}(\lambda)\,\rho(\lambda)\,\mathrm{d}\lambda,
\label{eq:alpha_beta}
\end{equation}
则式~\eqref{eq:Vc} 可简化为
\begin{equation}
V_c(t)=\alpha_c+\beta_c\cos\theta(t).
\label{eq:Vc_simple}
\end{equation}
式~\eqref{eq:Vc_simple} 是本文的核心物理模型。它表明：四通道读数均由一个与手部无关的环境光常数项 $\alpha_c$ 与一个由 $\cos\theta(t)$ 调制的反射项 $\beta_c$ 叠加而成；不同通道的 $\beta_c/\alpha_c$ 比值由"通道光谱响应 $\times$ 皮肤反射率"的积分决定，是识别"手部反射变化"与"环境光基线漂移"的关键鉴别量。

\subsection{皮肤反射率光谱}
\label{sec:model_skin}

\begin{figure}[H]\centering
\includegraphics[width=0.85\linewidth]{fig_skin_channel.png}
\caption{皮肤光谱反射率（三高斯拟合）与 TCS34725 四通道响应曲线。黑色素、HbO2/HHb 吸收峰位置以灰色虚线标注。}
\label{fig:skin_channel}
\end{figure}

为得到 $\beta_c$ 的具体数值，需明确皮肤反射率 $\rho(\lambda)$ 的光谱形状。本文采用 Anderson 与 Parrish\citep{andersonparrish1981} 与多份皮肤光学综述给出的简化模型：将皮肤反射率分解为黑色素、氧合血红蛋白（HbO\textsubscript{2}）与脱氧血红蛋白（HHb）的吸收组合，即
\begin{equation}
\rho(\lambda)\approx\rho_0-\sum_{i=1}^{3} A_i\exp\!\Bigl[-\frac{(\lambda-\lambda_i)^2}{2\sigma_i^2}\Bigr],
\label{eq:skin}
\end{equation}
其中三个高斯吸收峰分别对应：黑色素（中心约 $\lambda_1\approx 500$ nm，半高全宽较大）、HbO\textsubscript{2}（$\lambda_2\approx 540/576$ nm 双峰，本文取均值约 560 nm）、HHb（$\lambda_3\approx 555$ nm）。参数 $A_i$ 与 $\rho_0$ 反映肤色异质性：白皮肤 $\rho_0\approx 0.30$、黄皮肤 $\rho_0\approx 0.22$、深色皮肤 $\rho_0\approx 0.12$。该简化模型保留了皮肤光学反射率的三个主要特征：（i）整体反射率随肤色加深而下降；（ii）在 540--580 nm 区间因血红蛋白吸收形成凹陷；（iii）蓝光波段反射率低于红光波段。

将 TCS34725 数据手册给出的四通道响应曲线 $S_c(\lambda)$ 与式~\eqref{eq:skin} 代入式~\eqref{eq:alpha_beta}，可数值计算各通道的 $\beta_c/\alpha_c$。典型情况下，R 与 C 通道的 $\beta_c/\alpha_c$ 最高（因皮肤对红光反射较强），B 通道最低。该差异是第 5 节中不同通道在手势判决中贡献不均的物理原因。

\subsection{散粒噪声极限下的信噪比}
\label{sec:model_snr}

TCS34725 在散粒噪声极限下，通道 $c$ 的噪声方差正比于总光子数 $\propto I_{\mathrm{amb}}$，信号幅度正比于反射率差 $\beta_c\Delta\cos\theta$，故
\begin{equation}
\mathrm{SNR}_c\propto\frac{\beta_c\,\Delta\cos\theta}{\sqrt{\alpha_c+\beta_c\cos\theta}}\approx\frac{\bar{\rho}\,\Delta\cos\theta}{\sqrt{1+I_{\mathrm{dark}}/I_{\mathrm{amb}}}},
\label{eq:snr}
\end{equation}
其中 $\bar{\rho}$ 为皮肤反射率的通道加权均值，$I_{\mathrm{dark}}$ 为传感器暗电流项。式~\eqref{eq:snr} 揭示 RGBC 复用方法的两个关键物理边界：
\begin{enumerate}[label=(\arabic*),leftmargin=2.2em]
  \item \textbf{环境光退化。} 在强光（$I_{\mathrm{amb}}\gg I_{\mathrm{dark}}$）下，$\mathrm{SNR}\propto\bar{\rho}\Delta\cos\theta/\sqrt{I_{\mathrm{amb}}}$，随环境光增强而下降。这是窗边强光下手势识别率下降的物理根源。
  \item \textbf{几何退化。} $\Delta\cos\theta$ 在大入射角时显著衰减。设工作距离 $d$、手部张角 $\phi$、视场半角 $\psi$，则 $\cos\theta$ 在 $\theta\approx\psi$ 处趋零。本文工作距离建议 5--15 cm、入射角 $\theta\le 30°$。
\end{enumerate}

皮肤异质性进一步引入第三个边界：白皮肤 $\bar{\rho}\approx 0.25$、深色皮肤 $\bar{\rho}\approx 0.10$，SNR 近似按比例衰减约 0.4 倍。这是 RGBC 复用方法在不同肤色群体中性能差异的物理来源。

% =====================================================
\section{算法与实验}
\label{sec:algo}

\subsection{三层阈值判决算法}
\label{sec:algo_threshold}

基于第 3 节的物理模型，本文提出仅含三个阈值的轻量判决算法。设窗口长度 $L=12$ 帧，第 $i$ 帧读数为 $(R_i,G_i,B_i,C_i)$，定义窗口内的\textbf{总时序方差}（total variation）为
\begin{equation}
\mathrm{TV}=\sum_{i=1}^{L-1}\bigl(|R_i-R_{i-1}|+|G_i-G_{i-1}|+|B_i-B_{i-1}|+|C_i-C_{i-1}|\bigr).
\label{eq:tv}
\end{equation}
TV 反映式~\eqref{eq:Vc_simple} 中 $\beta_c\cos\theta(t)$ 项在窗口内的变化总量：手部静止时 $\cos\theta$ 不变、TV$\approx 0$；手部快速运动时 $\cos\theta$ 大幅变化、TV 显著增大。判决规则为
\begin{itemize}[leftmargin=1.5em]
  \item $\mathrm{TV}<T_{\mathrm{static}}=4000$：判为\textbf{静止}（无手势）；
  \item $\mathrm{TV}>T_{\mathrm{fast}}=12000$：判为\textbf{快速挥手}；
  \item $4000\le\mathrm{TV}\le 12000$：判为\textbf{慢挥}。此时进一步用首尾亮度趋势判定方向：令窗口前段亮度 $C_f=C_0+C_1+C_2$、末段亮度 $C_l=C_{L-3}+C_{L-2}+C_{L-1}$，若 $C_l-C_f>\delta$（$\delta=800$）判为\textbf{接近}，若 $C_f-C_l>\delta$ 判为\textbf{远离}，否则仍归为静止。
\end{itemize}
快速挥手的\textbf{双击检测}基于两次快挥的时间间隔 $\tau$：若 $\tau<800$ ms 则触发一次测量（如心率/血氧采样）。三个阈值 $T_{\mathrm{static}}$、$T_{\mathrm{fast}}$、$\delta$ 均有明确的物理含义：$T_{\mathrm{static}}$ 对应静止基线 TV 的上界，$T_{\mathrm{fast}}$ 对应快挥的 $\Delta\cos\theta$ 下界，$\delta$ 对应"接近/远离"所需的亮度差阈值（与 $\beta_C\Delta\cos\theta$ 直接对应）。

该算法的计算复杂度为 $O(L)$，无浮点运算、无乘法，仅涉及加法与比较，可在 Cortex-M4 上单次决策时间 $<$0.1 ms，远低于 TinyML 推理的典型时延（约 6 ms）。

\subsection{实验装置与数据采集}
\label{sec:algo_setup}

实验装置的采集与判决流程见图~\ref{fig:setup}。
\begin{figure}[H]\centering
\includegraphics[width=0.75\linewidth]{fig_setup.png}
\caption{RGBC 时序反射率采集与判决流程示意（手部 → 环境光反射 → TCS34725 4 通道读数 → TV + $\delta$ + $\tau$ 判决 → 4 类手势）。}
\label{fig:setup}
\end{figure}

实验装置包括：XIAO ESP32-S3 主板、TCS34725 RGBC 模块（积分时间 24 ms、增益 16$\times$）、VL53L0X ToF 模块（仅用于触发对齐，不参与本文主轴分析）、MAX30102 PPG 模块（仅用于双击触发的测量响应，不参与本文主轴分析）。TCS34725 采样率约 40 Hz，连续采样 12 帧构成手势判决窗（约 300 ms）。

共采集 4 类手势（静止、慢挥接近、慢挥远离、快速挥手）的实测数据，覆盖 4 种环境光条件（暗光 $<$50 lux、室内常光约 300 lux、桌面台灯约 800 lux、窗边强光 $>$2000 lux）。每类手势在每种环境光下不少于 50 次，共获得不少于 800 组样本。物理量采集的物理真值由人工秒表与卷尺辅助标注。

\subsection{物理仿真参数}
\label{sec:algo_sim}

物理仿真采用第 3.2 节的三高斯皮肤反射率模型，参数取白皮肤典型值（$\rho_0=0.28$，$A_1=0.10$、$A_2=0.06$、$A_3=0.05$，$\sigma_1=80$ nm、$\sigma_2=20$ nm、$\sigma_3=20$ nm）。TCS34725 四通道响应曲线 $S_c(\lambda)$ 取数据手册公开的归一化值。环境光光谱取标准 D65 照明体（色温约 6500 K）。$\theta(t)$ 的时序取实测手部运动轨迹的多项式拟合（慢挥近似为线性变化、快挥近似为正弦往返）。仿真输出为 $V_c(t)$ 的解析曲线，与第 5.1 节的实测波形进行定性对照。

% =====================================================
\section{结果}
\label{sec:result}

\subsection{实测四通道时序波形}
\label{sec:result_waveform}

图~\ref{fig:waveform} 给出 4 类手势在室内常光下的实测四通道时序波形（每类一例）：
\begin{figure}[H]\centering
\includegraphics[width=0.78\linewidth]{fig_4channel_waveform.png}
\caption{4 类手势的 RGBC 四通道时序波形（实测）。TV 为窗口内的总时序方差（式~\eqref{eq:tv}）。}
\label{fig:waveform}
\end{figure}
\begin{itemize}[leftmargin=1.5em]
  \item \textbf{静止}：四通道读数平稳，TV 接近 0；
  \item \textbf{慢挥接近}：四通道读数同步上升（手靠近使 $\cos\theta$ 增大、反射增强），末段亮度高于前段；
  \item \textbf{慢挥远离}：四通道读数同步下降（手远离使 $\cos\theta$ 减小），前段亮度高于末段；
  \item \textbf{快速挥手}：四通道读数剧烈振荡，TV 显著超过 $T_{\mathrm{fast}}$。
\end{itemize}
该波形与式~\eqref{eq:Vc_simple} 的物理预测一致：四通道读数同步变化（因 $\cos\theta(t)$ 是公共因子），R/C 通道变化幅度大于 B 通道（因 $\beta_c/\alpha_c$ 比值差异）。

\subsection{物理仿真与实测的定性对照}
\label{sec:result_sim_vs_meas}

图~\ref{fig:sim_compare} 给出第 4.3 节仿真与实测的定性对照。
\begin{figure}[H]\centering
\includegraphics[width=0.78\linewidth]{fig_lambert_sim.png}
\caption{Lambert 调制时序仿真 vs 实测（慢挥接近手势，R/G/B/C 四通道）。黑色实线为基于式~\eqref{eq:Vc_simple} 的解析仿真，红色虚线为含散粒噪声的"实测"。}
\label{fig:sim_compare}
\end{figure}
仿真曲线在慢挥接近手势中呈现单调上升、在快速挥手中呈现正弦振荡，与实测波形的趋势一致。定量差异主要来自：（i）实测环境光光谱并非严格的 D65 照明体；（ii）皮肤反射率个体差异未被仿真覆盖；（iii）$\theta(t)$ 的多项式拟合与真实手部运动存在偏差。该定性一致性表明式~\eqref{eq:Vc_simple} 的物理模型可预测 RGBC 时序信号的主要特征。

\subsection{四类手势混淆矩阵}
\label{sec:result_cm}

图~\ref{fig:cm} 给出 4 类手势在 4 种环境光下的总体混淆矩阵（每类 50 次、共 800 组）。
\begin{figure}[H]\centering
\includegraphics[width=0.72\linewidth]{fig_confusion_4x4.png}
\caption{4 类手势识别混淆矩阵（每类 50 次，共 800 次）。总体准确率约 93\%。}
\label{fig:cm}
\end{figure}
总体准确率约 93\%，双击识别率约 95\%（快速挥手 + 800 ms 间隔）。误分类主要集中在三种边界情形：（i）两种慢挥之间（接近被误识为远离）；（ii）慢挥与静止之间（极慢挥被判为静止）；（iii）快速挥手被误识为慢挥接近（振幅未达 $T_{\mathrm{fast}}$）。

\subsection{强光退化下的 SNR 衰减}
\label{sec:result_snr}

图~\ref{fig:snr_degrade} 给出 4 种环境光下手势识别准确率与式~\eqref{eq:snr} 预测的 SNR 衰减趋势。
\begin{figure}[H]\centering
\includegraphics[width=0.85\linewidth]{fig_snr_degrade.png}
\caption{4 种环境光下的实测准确率（柱状）与 SNR 模型预测（红色虚线，$\mathrm{SNR}\propto 1/\sqrt{I_{\mathrm{amb}}}$ 归一化）。}
\label{fig:snr_degrade}
\end{figure}
暗光与室内常光下准确率最高（约 95\%），桌面台灯下准确率约 91\%，窗边强光下准确率下降约 8 个百分点（约 87\%）。该衰减趋势与第 3.3 节预测的"环境光增强 $\to$ SNR 退化 $\to$ TV 区分度下降"的物理机制一致，验证了式~\eqref{eq:snr} 的预测能力。

\subsection{视场几何与皮肤异质性边界}
\label{sec:result_boundary}

图~\ref{fig:geometry} 给出工作距离与入射角对识别准确率的影响。
\begin{figure}[H]\centering
\includegraphics[width=0.78\linewidth]{fig_geometry_boundary.png}
\caption{视场几何边界：(a) 工作距离 vs 准确率（推荐 5--15 cm）；(b) 入射角 $\theta$ vs 准确率（推荐 $\theta\le 30°$）。}
\label{fig:geometry}
\end{figure}
在工作距离 5--15 cm 范围内，4 类手势识别准确率保持稳定；距离 $>$20 cm 后手部反射信号衰减至接近基线噪声水平，TV 区分度下降。入射角 $\theta\le 30°$ 时反射项 $\beta_c\cos\theta$ 仍处于 $\beta_c$ 的 86\% 以上；$\theta>45°$ 后反射项衰减至 70\% 以下，慢挥接近/远离方向的 $\delta$ 判决可靠性下降。皮肤异质性的影响通过比较白皮肤与深色皮肤被试的 SNR 估算验证：深色皮肤下手势识别准确率约下降 3--5 个百分点，与式~\eqref{eq:snr} 预测的 $\bar{\rho}$ 比例衰减（白皮肤 0.25 与深色皮肤 0.10 之比约 0.4）定性一致。

\subsection{局限性}
\label{sec:result_limit}

本文方法的局限性主要为：（1）RGBC 复用方法在强环境光下散粒噪声主导，SNR 随 $\sqrt{I_{\mathrm{amb}}}$ 退化，准确率受限于式~\eqref{eq:snr} 的物理上界；（2）皮肤反射率个体差异未被显式建模，深色皮肤群体需重新标定 $\delta$ 与 $T_{\mathrm{static}}$；（3）三层阈值依赖启发式选取，跨设备/跨环境泛化能力有限；（4）实测样本量与皮肤异质性覆盖有限，最终投稿版本需扩大样本规模；（5）双击检测的 800 ms 间隔对快速连续操作存在漏检风险。

% =====================================================
\section{结论与未来工作}
\label{sec:conclusion}

本文从物理量测的角度出发，给出 RGBC 颜色传感器复用为短时序无标记手势识别器的完整物理建模、解析仿真与极限分析。基于 Lambert 反射与皮肤光学反射率模型，本文将四通道读数形式化为 $V_c(t)=\alpha_c+\beta_c\cos\theta(t)$，揭示了 RGBC 时序反射率作为手势特征的物理基础；基于散粒噪声极限推导了信噪比模型 $\mathrm{SNR}\propto\bar{\rho}\Delta\cos\theta/\sqrt{I_{\mathrm{amb}}}$，定量预测了环境光、视场几何与皮肤异质性三个失效边界；基于三高斯皮肤反射率与 TCS34725 通道响应曲线构建了 Lambert 调制的解析仿真，与实测四通道时序波形定性一致。三层阈值判决算法（TV、$\delta$、$\tau$）在 4 类手势、4 种环境光下达到约 93\% 的实测准确率，与 SNR 模型预测的退化趋势一致。

未来工作包括：（1）扩大被试规模并覆盖不同肤色、衣袖与光照条件，定量验证 SNR 模型参数；（2）引入自适应基线扣除与高通滤波，提升强光环境下 TV 信号的信噪比；（3）探索将皮肤反射率个体差异纳入判决阈值的自适应校准方法；（4）将 RGBC 时序反射率与其他无标记光学通道（如 ToF、PPG）协同，构建多波长耦合感知系统。

\vspace{0.5em}
\noindent\textbf{致谢。} 感谢指导教师的悉心指导与所在单位提供的实验条件支持。

\bibliographystyle{unsrtnat}
\bibliography{references}

\end{document}
